% Include as a fragment in the patterns handbook; it deliberately has no chapter.
\section{Graph and DP casebook}

These cases practice a repeatable decision: identify what the answer depends
on, choose the smallest state that preserves that information, and only then
choose a traversal or data structure. The reasoning appears in prose and
short pseudocode, with links to tested local helpers where one matches the
problem contract. A helper is a reusable implementation component, not a
complete submission.

\subsection{Shortest route with a parity condition}
\label{case:state-bfs}
This is a synthetic graph example.

Suppose an undirected graph has unit-cost edges, and the task asks whether a
walk from $s$ to $t$ uses an even number of edges, or asks for the shortest
such walk. A plain BFS on vertices looks tempting: it gives the shortest
distance to every vertex. But knowing that a vertex is reachable does not
tell us whether it is reachable after an even or odd number of steps. If a
vertex is reached by paths of lengths two and three, those histories have
different eligibility for the destination condition. Keeping only one
distance per vertex throws away relevant information.

Use a product graph whose state is $(v,p)$, where $v$ is the current vertex
and $p\in\{0,1\}$ is the number of traversed edges modulo two. Each original
edge $(u,v)$ creates transitions $(u,p)\to(v,1-p)$ and
$(v,p)\to(u,1-p)$. The invariant is: BFS distance to $(v,p)$ is the fewest
edges in any walk from $s$ to $v$ whose length has parity $p$. Initialize
only $(s,0)$ at distance zero. Then a shortest even walk to $t$ is the
distance at $(t,0)$; an unreachable state remains absent. If actual
reconstruction is required, retain a parent state for each visited product
state and walk backward from $(t,0)$.

Pseudocode:
\begin{lstlisting}
dist[n][2] = UNREACHED
dist[s][0] = 0; queue.push((s, 0))
while queue not empty:
    (u, p) = queue.pop_front()
    for v in graph[u]:
        q = 1 - p
        if dist[v][q] is UNREACHED:
            dist[v][q] = dist[u][p] + 1
            parent[v][q] = (u, p)
            queue.push((v, q))
\end{lstlisting}

The expanded graph has $2n$ states and $O(m)$ transitions, so time and memory
stay $O(n+m)$. For ordinary reachability,
the local \verb@bfs(g,{source})@ helper is enough. For parity, represent
$(v,p)$ as \verb@2*v+p@ and construct the implicit transitions while
traversing, or build a graph with $2n$ vertices and call the helper. Its
returned parent array then supports \verb@restore_path@ on expanded state
IDs; decode each ID by division and remainder by two. A direct implementation
may store two distances per vertex instead.

Tiny trace: let the graph be the triangle $0$--$1$--$2$--$0$, with $s=0$.
The initial state is $(0,0)$. One BFS layer discovers $(1,1)$ and $(2,1)$.
The next layer discovers $(2,0)$ through 1 and $(1,0)$ through 2. Thus the
shortest even walk from 0 to 2 has length 2, namely $0,1,2$, even though the
shortest unrestricted walk has length 1. This illustrates why ``visited
vertex'' is not a sufficient marker.

The model changes if edges carry non-unit costs: parity still belongs in the
state, but BFS no longer minimizes cost. Use Dijkstra on the product graph
for nonnegative costs, or the appropriate weighted shortest-path method if
negative edges are allowed. If the condition is ``exactly $k$ edges,'' a
modulo-two state is insufficient; use a step dimension when $k$ is small, or
a method suited to the bound. If the question asks for a simple path, the
expanded walk can revisit an original vertex with a different parity, so the
walk-based model alone does not enforce simplicity.

Practice variant: find the shortest walk whose length is congruent to 1
modulo 3. Answer: use state $(v,p)$ with $p\in\{0,1,2\}$; every edge moves
to $(v,(p+1)\bmod 3)$; start at $(s,0)$ and read $(t,1)$. The product graph
has $3n$ states. This generalization works because the condition depends
only on the edge count modulo 3.

\subsection{A grid with moving hazards: CSES Monsters}
\label{case:hazards}

In this grid task the player and every monster can move one cell per turn;
the player must reach any boundary cell without ever occupying a cell with a
monster. A single BFS from the player is the obvious first attempt, but it
answers the wrong question. A cell that is reachable in three steps might
already be occupied by a monster at time two. Treating monsters as static
walls ignores the timing constraint; moving them in the same queue as the
player also mixes two different arrival-time questions.

The sufficient state has two parts conceptually: the player's cell and the
time at which the player arrives. Rather than include time in a large
time-expanded graph, first compute the earliest monster arrival time
$M[x]$ for every cell using one multi-source BFS from all monster cells.
Then compute a player BFS from $A$, entering a neighbor $y$ only if the
player's arrival time $P[x]+1$ is strictly less than $M[y]$. Equality is
unsafe because both arrive during the same turn. The invariant for the
second BFS is that each queued cell is reached by a valid player path and at
its shortest valid time; all such paths are safe because every earlier
arrival along that path passed the same strict test.

Steps:
\begin{enumerate}
\item Run multi-source grid BFS from every monster; walls are impassable.
Unreached cells have $M=\infty$.
\item Start player BFS at $A$ with $P[A]=0$. For each open neighbor $y$,
accept it only when it is unvisited and $P[x]+1<M[y]$.
\item Stop when a valid boundary cell is dequeued. Store its parent cell
and move direction at discovery time; reverse those moves to print a path.
If no boundary is reached, report failure.
\end{enumerate}
\columnbreak
\sourceimport{finals-patterns}{HazardEscape.h}
The tested \verb@finals_cases::escape_monsters(board)@ helper packages these
steps and returns an optional UDLR path; an empty optional means no escape.
It accepts the statement's \verb@.@, \verb@#@, \verb@A@, and \verb@M@
alphabet. Use the returned path directly when one valid route is enough.

This is $O(hw)$ time and space for an $h$ by $w$ board, since each cell is
processed a constant number of times. The local \verb@grid_bfs(board,sources,'#')@
helper computes ordinary grid distances with four-neighbor motion; call it
once for monsters and once for the player only if you can apply the safety
filter during player expansion. A plain call for player distances followed
by a check of only the final boundary is incorrect, because an unsafe
intermediate cell may have been traversed. For that constrained BFS, reuse
the helper's neighbor and parent pattern while adding the arrival test.
Flatten cells with \verb@id = row*w + col@ and decode with division and
remainder by \verb@w@. The statement permits up to $1000\times1000$ cells,
so the product and all stored IDs must fit the chosen integer type.

Tiny trace: consider rows \verb@#######@, \verb@#.....#@,
\verb@#M.A..#@, \verb@#.....#@, and \verb@#####.#@. The player starts at
row 2, column 3. The left neighbor has monster arrival time 1, so entering
it at time 1 fails the strict test. A shortest safe route goes down, right,
right, then down to the only open boundary cell: it takes 4 steps. The
monster at $(2,1)$ needs 6 steps to reach that exit at $(4,5)$, so the
player arrives safely. This
immediate comparison prevents the
mistake of declaring an escape merely because a boundary is geometrically
reachable.

Boundary details matter: if the player begins on a boundary, the zero-step
path is valid before any hazard can arrive, so check the start as an exit.
Unreachable monster distances behave as infinity; do not accidentally
interpret a default zero as an arrival at time zero. Monsters can be
initialized together because all spread at the same speed. If hazards have
different speeds, doors, or time-dependent routes, a single unweighted
multi-source BFS may no longer represent earliest arrival; use the matching
weighted or time-dependent model. If the player has a different speed, scale
times consistently or compare arrival times in common units.

Practice variant: hazards move twice as fast as the player, and all moves
remain on unit grid edges. Answer: monster BFS gives edge distance $d_M$;
the hazard time is $d_M/2$, so compare player arrival $d_P$ by cross
multiplying $2d_P<d_M$ (assuming positions change at half-turn granularity).
If turns are indivisible and the rules specify simultaneous discrete moves,
derive the exact timing convention from the statement before applying this
formula.

Source: the primary \href{https://cses.fi/problemset/task/1194/}{CSES Monsters statement} specifies simultaneous moves, adversarial knowledge of the path, and a boundary escape. The two-BFS method above is the state reduction used to enforce those rules.

\subsection{Digit completion as a remainder automaton: 2017 Enigma}
\label{case:enigma}

Enigma gives a fixed-length decimal pattern containing digits and question
marks, and asks for its smallest completion divisible by $N$, with no leading
zero. A tempting approach fills unknown digits with their smallest allowed values
and tries to repair divisibility later. This fails because a locally small digit can leave a suffix remainder
that no legal completion can fix. Enumerating all replacements is
exponential in the number of question marks. The key observation is that a
prefix's entire numeric value is unnecessary: future divisibility depends
only on its remainder modulo $N$.

Define $F[i][r]$ to mean: starting before pattern position $i$ with current
prefix remainder $r$, there exists a legal suffix completion that ends with
remainder zero. The recurrence tries each allowed digit $d$ and sets
$F[i][r]$ if $F[i+1][(10r+d)\bmod N]$ is true. Base case:
$F[L][0]=\text{true}$ and $F[L][r]=\text{false}$ for $r\ne0$. At position
zero, disallow digit zero. To reconstruct the lexicographically smallest
number, start with remainder zero and choose the smallest allowed digit
whose next state is feasible in the next row. This greedy choice is safe
because the table certifies that some valid suffix remains.

The local \verb@finals_dp::complete_digits(pattern,modulus)@ implements this
API directly: it returns an optional smallest string, forbids a leading zero
by default (including a one-character zero), and runs in $O(10LN)$ time and
$O(LN)$ memory. Its boolean table is filled backward; reconstruction walks
forward. For a larger modulus, this memory bound may be the limiting factor,
so check constraints before allocating $(L+1)N$ entries. The official 2017
statement has $L\le1000$ and $N\le1000$, making the table the intended
shape.

Tiny trace: pattern \verb@?2?@ and modulus 7. The smallest legal first digit
is 1. After choosing 1 the prefix remainder is 1. At position 1 the fixed
digit 2 changes it to $(10\cdot1+2)\bmod7=5$. The last digit $d$ must satisfy
$(10\cdot5+d)\bmod7=(1+d)\bmod7=0$, so $d=6$. The answer is 126. Trying
first digit 0 is forbidden; choosing 1 is justified by suffix feasibility,
not by guessing.

The official statement allows the initial character to be a question mark
but requires the completed number not to start with zero; its samples also
include an impossible pattern. Keep ``no completion'' distinct from a valid
answer containing zero. A common implementation bug is to compute the
forward remainder correctly but use a suffix table with the wrong meaning;
state the table meaning before writing the recurrence. Another is to use
\verb@10*r+d@ in a narrow integer type when bounds could overflow; the local
helper widens before multiplication.

Modification: if the task asks for the number of completions rather than
the smallest one, replace boolean feasibility with counts and sum over
allowed digits. The local helper does not do that; moreover, counts may need
a modulus or arbitrary precision. If the task asks for the smallest number
at most a bound, add a tight-to-bound state as well as the remainder. The
remainder alone is no longer sufficient because the bound depends on the
prefix relation to the limit.

Practice variant: complete \verb@?3@ to a multiple of 4. Answer: no
completion exists, because every multiple of 4 is even while the fixed final
digit is 3. The feasibility table reports that directly; there is no
candidate digit or retraction to apply.

Source: the primary \href{https://maratona.ime.usp.br/hist/2017/resultados/contest.pdf}{2017 Enigma statement} gives the pattern, divisibility, leading-zero rules, and bounds. The remainder-DP explanation matches the tested local helper.

\subsection{Weighted dominance DP with strict ties: 2017 Fundraising}
\label{case:fundraising}

Each candidate has beauty $B$, fortune $F$, and donation $D$. Two people
can coexist only when they are equal in both attributes or one is strictly
larger in both. We seek the maximum donation sum. Sorting by one coordinate
and taking a weighted LIS is close, but a naive implementation can admit
illegal ties: processing points with equal fortune one at a time may let a
later point use an earlier one, even though equal fortune does not satisfy
the strict increase rule. A prefix sum Fenwick tree is also the wrong query:
we need the best predecessor value among compatible points, not the sum of
all their values.

First merge identical $(B,F)$ pairs and sum their donations. It is safe and
necessary: equal-coordinate people do not argue and can all be invited, so
they act as one point of combined weight. Then sort by fortune ascending and,
for equal fortune, beauty descending. For a point $i$, its DP value is
$D_i+\max(0,\text{best donation sum ending at a point with }F<F_i,B<B_i)$.
The zero option starts a new chain. The best valid guest set forms a chain
under strict increase in both coordinates: any pair in the chain can
coexist, and the statement's compatibility condition excludes incomparable
unequal pairs.

Compress beauty values to zero-based ranks $0..K-1$. A prefix-maximum
Fenwick tree can answer the maximum DP among beauty ranks strictly below the current rank,
then update the current rank with the candidate DP. If only the generic local
\verb@SegTree<T,Op>@ is available, use max with identity zero (donations are
positive), query half-open ranks $[0,rank)$, and point-set the rank to the
maximum of its current value and the new DP. Do not use a range sum or the
ordinary additive \verb@Fenwick@ helper: neither represents prefix maximum.
Also, point-set must preserve the previous maximum if another valid chain
ends at the same rank.

There are two correct tie strategies. One is to sort by fortune ascending
and beauty descending; a sequential strict-prefix query cannot see equal
beauty, and descending beauty prevents equal-fortune points from increasing
in beauty. Alternatively, process all equal-fortune points as a batch: query
all of them against the tree before applying any updates.
Batching is often easier to reason about and remains correct with another
tie sort. Use one strategy deliberately, not half of each.

The tested \verb@finals_cases::max_donation_chain(donors)@ helper in this
booklet applies the batched strict-coordinate method. Its \verb@Donor@
fields are \verb@wealth@, \verb@beauty@, and \verb@donation@; nonnegative
donations are supported, and identical coordinate pairs are combined.
\columnbreak
\sourceimport{finals-patterns}{Fundraising.h}
Use it when the input matches that contract; keep the recurrence below as
the explanation for why the data-structure query is correct.

Pseudocode for the editorial ordering:
\begin{lstlisting}
merge donations for each identical (beauty, fortune)
sort by (fortune ascending, beauty descending)
compress all beauty coordinates
for (B, F, D) in sorted points:
    r = rank(B)
    best = prefix_max(0, r)          // zero-based half-open ranks [0,r)
    dp = best + D
    update_prefix_max(r, dp)
answer = maximum dp (or tree-wide maximum)
\end{lstlisting}

The local segment tree uses zero-based ranks and half-open ranges, so query
$[0,r)$. A Fenwick implementation may instead use one-based update indices;
then query the first $r$ entries. The local Fenwick \verb@prefix(r)@ API
also means $[0,r)$. Define the index convention and prefix API together so
the exclusion of equal beauty is visible in code. Complexity is $O(N\log N)$ time and $O(N)$
space, with $N\le10^5$; donation sums require 64-bit integers because each
donation may be as large as $10^9$.

Tiny trace: points $(B,F,D)=(1,1,4),(2,1,7),(3,2,5)$. Sort at fortune 1 by
beauty descending: $(2,1,7)$ then $(1,1,4)$. Each starts its own chain because
neither can see an earlier strictly smaller beauty at the same fortune. At
$(3,2,5)$, the prefix query sees both prior values and returns 7, giving 12.
A fortune-ascending, beauty-ascending sort would process $(1,1,4)$ before
$(2,1,7)$ and incorrectly chain them for 11, then report 16. Merging equal
points matters too: two guests at $(2,1)$ with donations 3 and 8 may both be
invited, so their grouped weight is 11.

The maximum query uses identity zero because the statement's donations are
positive and starting a new chain is allowed. If negative donations were
permitted but the guest list had to be nonempty, initialize the final answer
to negative infinity; for each point, set $dp=D+\max(0,\text{best
predecessor})$, so zero starts a new nonempty chain at that point. Do not
initialize the answer to zero, which would represent an invalid empty guest
list when every donation is negative. Also, merging identical coordinate
pairs by summing weights relies on the original positive donations: with
negative donations, one could invite only some people in that group, so
retain their individual choices. If the problem changed to allow equality in beauty,
query through the current rank; if it allowed equality in fortune, equal-$F$
points could interact and the descending-tie trick would no longer enforce
strictness in that coordinate. For reconstruction, store the chosen
predecessor point whenever a new maximum is written; ties can use any
predecessor meeting the constraints.

Practice variant: require beauty strictly increasing but permit equal
fortune. Answer: keep fortune ascending, and query beauty ranks strictly
below as before; equal-fortune points may now chain, so process them in
ascending beauty order and allow each update to affect the next point. This
change is valid only because the modified rule explicitly permits equal
fortune.

Source: the primary \href{https://maratona.ime.usp.br/hist/2017/resultados/contest.pdf}{2017 Fundraising statement} defines compatibility, donation bounds, and input size. The casebook's recurrence and tie strategies are checked against the tested local helper; see \verb@handbooks/research/graph-dp-pattern-sources.md@ for source-access notes.
